\ncuAbstractPage{\ncuENtitle}{英文摘要}{ABSTRACT}
Temporal change and class imbalance often occur together in classification. Historical data provide more samples but may no longer reflect recent conditions, whereas recent data are closer to the prediction period but contain fewer minority-class cases. This study examines how models built from historical (Old) and recent (New) data can be used for time-stamped imbalanced classification, and evaluates the effects of ensemble learning and feature selection.

The empirical case uses public annual financial data of U.S. companies from 1999 to 2018. A next-year bankruptcy event label is reconstructed from the last observed year of each failed company, so that positive cases account for only 0.77\% of all records. Experiment 1 compares Old, New, and combined retraining (Retrain) for four classifiers under the same year configuration, and then builds a six-model Old/New XGBoost pool to compare static ensembles, dynamic classifier selection, and dynamic ensemble selection. Experiment 2 adds four feature selection methods with two retention levels before sampling. Ensemble members and decision thresholds are selected on validation data only.

The advantage of recent data depends on the classifier. Within the fixed pool, the New-only ensemble (New3) achieves the highest mean AUC (0.8601), and the validation-selected static ensemble, which keeps two members, achieves the highest F1-score (0.1795). Equal weighting of all six models and dynamic selection do not improve performance. Feature selection does not yield consistent gains across methods and metrics. Tuned single classifiers perform comparably to the ensembles, so the evidence supports careful selection of ensemble members rather than a general superiority of ensembles. When the companies' final status is used as the label instead, the ranking of methods changes, indicating that the label definition affects the conclusions.

\noindent\textbf{Keywords:} class imbalance, temporal data, ensemble learning, dynamic selection, feature selection, bankruptcy prediction
